\begin{solution}{109}
\begin{subsolution}
在相对于 $\pi$ 介子静止的参照系 $S'$（$\pi$ 介子的质心系）中，由动量守恒有
\begin{equation}
p_{\mu}' + p_{\nu}' = 0
\label{eq:1.s109}
\end{equation}
式中 $p_{\mu}'$ 和 $p_{\nu}'$ 分别是 $\mu$ 子和 $\mu$ 型反中微子 $\overline{\nu}_{\mu}$ 在 $S'$ 系中的动量
\begin{equation}
p_{\mu}' = \frac{m_{\mu} u_{\mu}'}{\sqrt{1 - \left(\frac{u_{\mu}'}{c}\right)^{2}}}
\label{eq:2.s109}
\end{equation}
由能量守恒有
\[
E_{\mu}' + E_{\nu}' = m_{\pi} c^{2}
\]
式中 $E_{\mu}'$ 和 $E_{\nu}'$ 分别是 $\mu$ 子和 $\mu$ 型反中微子 $\overline{\nu}_{\mu}$ 在 $S'$ 系中的能量
\begin{equation}
E_{\mu}' + E_{\nu}' = m_{\pi} c^{2}
\label{eq:3.s109}
\end{equation}
\begin{equation}
E_{\mu}' = \frac{m_{\mu} c^{2}}{\sqrt{1 - \left(\frac{u_{\mu}'}{c}\right)^{2}}},\quad E_{\nu}' = (-p_{\nu}') c
\label{eq:4.s109}
\end{equation}
将\eqref{eq:2.s109}和\eqref{eq:4.s109}式代入\eqref{eq:1.s109}和\eqref{eq:3.s109}式得
\[
\frac{1 + \frac{u_{\mu}'}{c}}{1 - \frac{u_{\mu}'}{c}} = \frac{m_{\pi}^{2}}{m_{\mu}^{2}}
\]
即
\begin{equation}
u_{\mu}' = \frac{m_{\pi}^{2} - m_{\mu}^{2}}{m_{\pi}^{2} + m_{\mu}^{2}} c = 0.2714c
\label{eq:5.s109}
\end{equation}

[另一解法:
在相对于 $\pi$ 介子静止的参照系 $S'$（$\pi$ 介子的质心系）中，由动量守恒有
\[
p_{\mu}' + p_{v}' = 0
\]
式中 $p_{\mu}'$ 和 $p_{v}'$ 分别是 $\mu$ 子和 $\mu$ 型反中微子 $\overline{v}_{\mu}$ 在 $S'$ 系中的动量；由能量守恒有
\[
(E_{\mathrm{k}\mu}' + m_{\mu} c^{2}) + E_{\nu}' = m_{\pi} c^{2}
\]
式中 $E_{\mathrm{k}\mu}'$ 和 $E_{\nu}'$ 分别是在 $S'$ 系中 $\mu$ 子和 $\mu$ 型反中微子 $\overline{\nu}_{\mu}$ 的动能。自由粒子的能量与动量满足
\[
E_{\mu}'^{2} = (p_{\mu}' c)^{2} + (m_{\mu} c^{2})^{2},\quad E_{\nu}' = p_{\nu}' c
\]
利用 $E_{\mu}' = \gamma_{\mu}' m_{\mu} c^{2}$，能量守恒式即
\[
E_{\nu}' = m_{\pi} c^{2} - \gamma_{\mu}' m_{\mu} c^{2}
\]
式中
\[
\gamma_{\mu}' = \frac{1}{\sqrt{1 - \frac{u_{\mu}'^{2}}{c^{2}}}}
\]
而 $u_{\mu}'$ 是 $\mu$ 子在 $S'$ 系中的运行速度。利用动量守恒和能量动量关系得
\[
E_{\nu}' = p_{\nu}' c = p_{\mu}' c = \sqrt{E_{\mu}'^{2} - (m_{\mu} c^{2})^{2}} = m_{\pi} c^{2} - \gamma_{\mu}' m_{\mu} c^{2} = (m_{\pi} - m_{\mu}) c^{2} - m_{\mu} c^{2} (\gamma_{\mu}' - 1)
\]
上式两边平方后解得
\[
\gamma_{\mu}' = 1 + \frac{(m_{\pi} - m_{\mu})^{2}}{2 m_{\pi} m_{\mu}} = 1.0390
\]
由此可导出 $\mu$ 子在 $S'$ 系中的速度为:
\[
u_{\mu}' = c \sqrt{1 - \frac{1}{\gamma_{\mu}'^{2}}} = 0.2714c
\]
]

假设在实验室系 $S$ 中，$\pi$ 介子以及衰变后的 $\mu$ 子的飞行方向之间的夹角为 $\theta$，这一夹角在 $S'$ 系中对应于夹角 $\theta'$。按相对论速度变换规则，在实验室系 $S$ 中，$\mu$ 子的飞行速度的大小 $u_{\mu}$ 满足
\[
u_{\mu} = \frac{\sqrt{u_{\mu}'^{2} + v^{2} + 2 u_{\mu}' v \cos\theta' - \left(\frac{u_{\mu}' v}{c}\right)^{2} \sin^{2}\theta'}}{1 + \frac{u_{\mu}' v}{c^{2}} \cos\theta'}
\]
此即
\[
u_{\mu} = c \sqrt{1 - \frac{\left(1 - \frac{u_{\mu}'^{2}}{c^{2}}\right)\left(1 - \frac{v^{2}}{c^{2}}\right)}{\left(1 + \frac{u_{\mu}' v}{c^{2}} \cos\theta'\right)^{2}}}
\]
由上式可知，从 $\pi$ 介子衰变得到的 $\mu$ 子的最大速率 $u_{\max}$ 和最小速率 $u_{\min}$ 分别出现在 $\theta' = 0$ 和 $\theta' = \pi$ 的情形，正好分别相应于 $\theta = 0$ 和 $\theta = \pi$ 的情形。于是由
\[
u_{\mu} \cos\theta = \frac{u_{\mu}' \cos\theta' + v}{1 + \frac{u_{\mu}' v}{c^{2}} \cos\theta'}
\]
可知，从 $\pi$ 介子衰变得到的 $\mu$ 子的最小速率 $u_{\min}$ 和最大速率 $u_{\max}$ 分别为
\begin{equation}
u_{\mu\min} = \frac{-u_{\mu}' + v}{1 + \frac{(-u_{\mu}') v}{c^{2}}} = 0.940c = 2.82 \times 10^{8}\,\mathrm{m/s}
\label{eq:6.s109}
\end{equation}
\begin{equation}
u_{\mu\max} = \frac{u_{\mu}' + v}{1 + \frac{u_{\mu}' v}{c^{2}}} = 0.980c = 2.94 \times 10^{8}\,\mathrm{m/s}
\label{eq:7.s109}
\end{equation}
\end{subsolution}

\begin{subsolution}
根据放射性衰变规律，$\mu$ 子的衰变率 $\rho$ 满足
\begin{equation}
\rho = \frac{N_{t}}{N_{0}} = 2^{- \frac{\Delta t}{T_{1/2}}}
\label{eq:8.s109}
\end{equation}
式中，$N_{0}$ 为初始 $\mu$ 子数，$N_{t}$ 为 $t$ 时刻的 $\mu$ 子数，$\Delta t$ 为 $\mu$ 子走完 $h = 10000\,\mathrm{m}$ 路程所需的时间间隔。$\Delta t$ 在 $u_{\mu} = u_{\mu\max}$ 情形下最短，从而在地面附近出现的概率最大，即
\begin{equation}
(\Delta t)_{\min} = \frac{h}{u_{\mu\max}} = 3.40 \times 10^{-5}\,\mathrm{s}
\label{eq:9.s109}
\end{equation}
由于时间膨胀效应，其半衰期在地面参照系中为
\begin{equation}
T_{1/2} = \frac{T_{1/2}^{(0)}}{\sqrt{1 - \frac{u_{\mu}^{2}}{c^{2}}}} = 0.765 \times 10^{-5}\,\mathrm{s}
\label{eq:10.s109}
\end{equation}
从而，在地面上观测到这个 $\mu$ 子的概率为
\begin{equation}
\rho_{\max} = 2^{- \frac{(\Delta t)_{\min}}{T_{1/2}}} = 4.59\%
\label{eq:11.s109}
\end{equation}

[另一解法:
根据放射性衰变规律，$\mu$ 子的衰变率 $\rho$ 满足
\[
\rho = \frac{N_{t}}{N_{0}} = \exp\left(- \frac{\Delta t}{\tau}\right)
\]
式中，$\tau$ 为 $\mu$ 子的平均寿命。$\Delta t$ 为 $\mu$ 子走完 $10000\,\mathrm{m}$ 路程所需的时间间隔，在 $u_{\mu} = u_{\mu\max}$ 情形下最短，因而在地面附近出现的概率最大，即
\[
(\Delta t)_{\min} = \frac{h}{u_{\mu\max}} = 3.40 \times 10^{-5}\,\mathrm{s}
\]
静止的 $\mu$ 子的平均寿命 $\tau_{0}$ 可从静止的 $\mu$ 子半衰期 $T_{1/2}^{(0)}$ 求得
\[
\tau_{0} = \frac{T_{1/2}^{(0)}}{\ln 2} = 2.197 \times 10^{-6}\,\mathrm{s}
\]
由于时间膨胀效应，其平均寿命在地面参照系下
\[
\tau = \frac{\tau_{0}}{\sqrt{1 - \frac{u_{\mu}^{2}}{c^{2}}}} = 1.10 \times 10^{-5}\,\mathrm{s} \;(\text{或 } 1.104 \times 10^{-5}\,\mathrm{s})
\]
从而，在地面上观测到这个 $\mu$ 子的概率为
\[
\rho_{\max} = \exp\left(- \frac{(\Delta t)_{\min}}{\tau}\right) = 4.55\% \; \text{或} \; 4.59\%
\]
]
\end{subsolution}
\end{solution}